\scanpage{39}

au voisinage de $s$. Donc, en ce voisinage

\[
u'=\underset{\substack{\Vert\\-b-u}}{\bar u}+b=\deleted{(\operatorname{Tr}(u)1-u)}-u=u
\]

et donc on aurait $\Theta=\mathrm{id}$ au voisinage de $s$, contrairement à l'hypothèse.

Donc on conclut $\alpha_s\in M_s$, mais alors, comme $b_s\notin M_s$, \deleted{\unsure{$b_s\in O_s^*$}} on aurait $\alpha_s-b_s\notin M_s$ i.e. $\alpha_s-b_s\in\underline O_s^*$, d'où (comme $\alpha_s(\alpha_s-b_s)=0$) $\alpha_s=0$, donc $\alpha$ nul au voisinage de $s$, cqfd.

\uline{NB} En car. $2$, dans le cas non étale, l'exemple typique est $A=\underline O[u]/(u^2)$. Les autom. de $A$ sont donnés \inserted{(indép. de la car.)} par

\[
\Theta_{\lambda,\alpha}:u\longmapsto\lambda u+\alpha\qquad\lambda\in\Gamma O^*,\ \alpha\in\Gamma\underline O
\]

avec

\[
\left\{\begin{aligned}2\alpha&=0\\\alpha^2&=0\end{aligned}\right.
\]

ce qui, en car. $2$, se réduit à

\[
\alpha^2=0.
\]

Ceux qui induisent $-\mathrm{id}_{A/O.1}$ sont ceux de la forme

\[
u\longmapsto-u+\alpha\qquad(\text{avec }2\alpha=0,\ \text{\uuline{$\alpha^2=0$}}),
\]

donc en car. $2$ \unsure{ceux} \uline{de la forme} $u\longmapsto u+\alpha$ avec $\alpha^2=0$.$^*$

% slanted
\margin{Dire qu'en car $2$, \deleted{\unsure{les} \unsure{tot.}} ramifié ($\operatorname{Tr}=0$ \unsure{i.e.} $\sigma=\unsure{0|u}$) les symétries $\sigma_D$ sont l'\uline{identité}\,!}

% slanted
\margin{Il n'y a pas que l'identité i.e. pour que $\sigma_D$ puisse être l'identité sur $D$, $-\mathrm{id}$ sur $E/D$ (en car. $2$, les \unsure{car.} \unsure{réunies})}

Soit $(E,\Delta,Q)$ une str. de similitude str. fid, et $D\subset E$ droite tot\textsuperscript{t} non isotrope, d'où l'autom $\sigma_D\in\Gamma\underline{\operatorname{End}}(E)^*$. Le fait que $\sigma_D$ soit un autom de la structure de similitude signifie

% archive sheet number: 38
